\documentclass[11pt]{article}
\usepackage[T1]{fontenc}
\usepackage{amsmath,amssymb,amsthm}
\usepackage[margin=1in]{geometry}
\usepackage{booktabs}
\usepackage{microtype}
\usepackage[hidelinks,pdfusetitle]{hyperref}

\newtheorem{theorem}{Theorem}
\newtheorem{lemma}[theorem]{Lemma}
\newcommand{\Aut}{\operatorname{Aut}}
\newcommand{\Orb}{\operatorname{Orb}}
\newcommand{\Stab}{\operatorname{Stab}}

\title{All Minimum Bondage Sets of the Truncated-Octahedral Graph, up to Automorphism}
\author{Xiaojian Zeng\\
PBC School of Finance, Tsinghua University, 100190\\
\texttt{xiaojianzeng@foxmail.com}}
\date{}

\begin{document}
\maketitle

\begin{abstract}
Let $T$ be the $24$-vertex truncated-octahedral graph in the recent
construction of Srivastava.  That construction establishes
$\gamma(T)=8$ and bondage number $b(T)=5$, and exhibits one five-edge
bondage set.  We classify all of them.  Exactly $276$ five-edge sets $F$
satisfy $\gamma(T-F)>\gamma(T)$.  The full abstract automorphism group of
$T$ has order $48$, and the $276$ sets form eight orbits, of sizes
$24,48,12,48,48,24,48,24$.  The corresponding stabilizer orders are
$2,1,4,1,1,2,1,2$, and $\gamma(T-F)=9$ in every orbit.  In particular,
every minimum bondage set contains the full three-edge star at a vertex.
The classification is certified by exhaustive integer-mask enumeration;
a separate automorphism backtracker independently reproduces the family
and orbit data.
\end{abstract}

\section{Statement and context}

A vertex set $D$ in a finite simple graph $G$ is \emph{dominating} if
every vertex outside $D$ has a neighbor in $D$, and $\gamma(G)$ is the
least cardinality of such a set.  Fink, Jacobson, Kinch, and Roberts
introduced the bondage number
\[
 b(G)=\min\{|F|:F\subseteq E(G),\ \gamma(G-F)>\gamma(G)\}
\]
as a measure of the sensitivity of domination to edge deletion
\cite{Fink1990Bondage}.

Srivastava proved that the truncated-octahedral graph $T$ below has
$\gamma(T)=8$ and $b(T)=5$ \cite{Srivastava2026Truncated}.  The proof
enumerates the $1{,}767$ dominating eight-sets, checks all $58{,}905$
four-edge deletions, and gives one five-edge witness.  It does not count
all five-edge bondage sets or classify them under graph automorphisms.
The result here supplies exactly that finite classification; it makes no
claim of a broader bondage theorem or of literature-wide priority.

For a textual specification, let $V(T)=\{0,1,\ldots,23\}$ and
\begin{align*}
E(T)=\{& (0,1),(0,2),(0,3),(1,4),(1,5),(2,6),(2,7),(3,4),(3,8),\\
       &(4,9),(5,10),(5,11),(6,11),(6,12),(7,12),(7,13),(8,13),(8,14),\\
       &(9,15),(9,16),(10,16),(10,17),(11,17),(12,18),(13,19),\\
       &(14,15),(14,19),(15,20),(16,20),(17,21),(18,21),(18,22),\\
       &(19,22),(20,23),(21,23),(22,23)\}.
\end{align*}
All pairs denote unordered edges.  Put
\[
 \mathcal M(T)=\{F\subseteq E(T): |F|=5,
              \ \gamma(T-F)>\gamma(T)\}.
\]
The full abstract group $\Aut(T)$ acts on $\mathcal M(T)$.  Within an
orbit, an edge is written with its smaller endpoint first, the five
edges are sorted lexicographically, and the least resulting tuple is
called the canonical representative.

\begin{theorem}\label{thm:classification}
The family $\mathcal M(T)$ has cardinality $276$, and $\Aut(T)$ has
order $48$.  There are exactly eight orbits, whose canonical
representatives and invariants are listed in Table~\ref{tab:orbits}.
For every $F\in\mathcal M(T)$ one has $\gamma(T-F)=9$, and $F$ contains
all three edges incident with some vertex.
\end{theorem}

\begin{table}[ht]
\centering
\small
\begin{tabular}{@{}llllr r r@{}}
\toprule
&\multicolumn{4}{c}{canonical representative $F_i$}
&$|\Orb(F_i)|$&$|\Stab(F_i)|$\\
\cmidrule(lr){2-5}
$i$&\multicolumn{4}{l}{five unordered edges}& &\\
\midrule
1&$(0,1),(0,2),(0,3)$&$(1,4)$&$(1,5)$&&24&2\\
2&$(0,1),(0,2),(0,3)$&$(1,4)$&$(4,9)$&&48&1\\
3&$(0,1),(0,2),(0,3)$&$(2,6)$&$(2,7)$&&12&4\\
4&$(0,1),(0,2),(0,3)$&$(2,6)$&$(8,14)$&&48&1\\
5&$(0,1),(0,2),(0,3)$&$(5,10)$&$(5,11)$&&48&1\\
6&$(0,1),(0,2),(0,3)$&$(5,10)$&$(8,14)$&&24&2\\
7&$(0,1),(0,2),(0,3)$&$(6,11)$&$(6,12)$&&48&1\\
8&$(0,1),(0,2),(0,3)$&$(6,11)$&$(7,13)$&&24&2\\
\bottomrule
\end{tabular}
\caption{The eight orbits in $\mathcal M(T)$.  In every row
$\gamma(T-F_i)=9$.}
\label{tab:orbits}
\end{table}

\section{The exact finite certificate}

The proof uses only finite iteration and integer bit operations.  We
give the criterion underlying the enumeration before reporting its
output.

\begin{lemma}[bundle criterion]\label{lem:bundle}
Let $\mathcal D_8$ be the family of dominating eight-sets of $T$.  For
$D\in\mathcal D_8$ and $x\notin D$, define
\[
 B(D,x)=\{xy\in E(T):y\in D\}.
\]
For $F\subseteq E(T)$, the set $D$ fails to dominate $T-F$ if and only
if $B(D,x)\subseteq F$ for some $x\notin D$.  Consequently, a five-edge
set $F$ is in $\mathcal M(T)$ if and only if
\[
 \forall D\in\mathcal D_8\quad
 \exists x\notin D\quad B(D,x)\subseteq F.
 \tag{1}\label{eq:criterion}
\]
\end{lemma}

\begin{proof}
A vertex of $D$ remains dominated by itself after any edge deletion.
An outside vertex $x$ ceases to be dominated by $D$ precisely when all
edges from $x$ to $D$ have been deleted, which is the stated bundle
condition.  Since $\gamma(T)=8$, an edge-deleted subgraph has domination
number greater than eight exactly when none of the sets in
$\mathcal D_8$ survives.  Any set that dominates $T-F$ also dominates
the supergraph $T$, so no other eight-set has been omitted.
\end{proof}

The verifier encodes every closed neighborhood by a $24$-bit integer.
It first checks directly that the displayed edge list is simple,
connected, and $3$-regular.  It enumerates every vertex subset of sizes
$0$ through $8$ and obtains the numbers of dominating sets
\[
 (0,0,0,0,0,0,0,0,1767).
 \tag{2}\label{eq:domcounts}
\]
Thus it independently recovers $\gamma(T)=8$ and $|\mathcal D_8|=1767$.
It then iterates through all
\(\binom{36}{5}=376{,}992\) five-edge masks and evaluates
\eqref{eq:criterion} exactly.  Precisely $276$ masks pass.  For each of
the other $376{,}716$ masks, the exact loop finds a
surviving dominating eight-set; every passing mask is replayed against
all $1{,}767$ sets.  This proves both inclusions in the reported family
census.

The first implementation enumerates self-isomorphisms with NetworkX
$3.4.2$ \cite{Hagberg2008NetworkX}, whose graph matcher implements the
VF2 approach \cite{Cordella2004VF2}.  Every returned map is separately
checked for bijectivity and for equality of the complete image edge set
with $E(T)$.  It returns $48$ maps.  The following three permutations,
in disjoint-cycle notation, generate all $48$ of them:
\begin{align*}
g_1={}&(1\ 3)(5\ 8)(6\ 7)(10\ 14)(11\ 13)(15\ 16)(17\ 19)(21\ 22),\\
g_2={}&(0\ 1)(2\ 5)(3\ 4)(6\ 11)(7\ 10)(8\ 9)(12\ 17)(13\ 16)\\
     &\qquad(14\ 15)(18\ 21)(19\ 20)(22\ 23),\\
g_3={}&(0\ 2)(1\ 6)(3\ 7)(4\ 12)(5\ 11)(8\ 13)(9\ 18)(10\ 17)\\
     &\qquad(14\ 19)(15\ 22)(16\ 21)(20\ 23).
\end{align*}

Completeness of the group computation was checked by a separately
implemented backtracker, without a graph-isomorphism library.  It tries
all $24$ images of vertex $0$ and all six bijections from its three
neighbors to the neighbors of the image.  It then exhaustively extends
each partial map, pruning only on adjacency and graph-distance
conditions that every automorphism must preserve.  Every automorphism
determines one of these initial seeds, so this search covers the full
abstract automorphism group.  Each completed map is checked against all
$36$ edges.  The independent search also returns exactly $48$ maps.

Acting with these maps on the $276$ accepted masks yields the eight
disjoint orbits of Table~\ref{tab:orbits}.  Their sizes sum to $276$, and
each row satisfies
\[
 |\Orb(F_i)|\,|\Stab(F_i)|=48.
\]
Lexicographic minimization over every orbit gives exactly the displayed
representatives.  This proves the orbit count, representatives, and
stabilizers.  Since each representative contains the three-edge star
$\{(0,1),(0,2),(0,3)\}$, automorphism invariance proves the final
structural assertion of Theorem~\ref{thm:classification}.

It remains to certify the post-deletion domination numbers.  By
Lemma~\ref{lem:bundle}, each accepted mask has domination number at least
nine.  For the representatives $F_1,\ldots,F_8$, respectively, the
following sets dominate after the indicated five edges are deleted:
\begin{align*}
W_1&=\{0,1,2,3,6,10,15,19,21\},\\
W_2&=\{0,1,2,3,6,9,16,19,21\},\\
W_3&=\{0,1,2,3,4,12,17,19,20\},\\
W_4=W_5=W_6=W_7=W_8
   &=\{0,1,2,3,4,12,17,19,20\}.
\end{align*}
The verifier checks the union of the nine closed-neighborhood masks in
each case against the all-vertices mask.  Hence
$\gamma(T-F_i)\leq 9$, and the lower bound gives equality.
Automorphisms carry these witnesses throughout their orbits, completing
the proof of Theorem~\ref{thm:classification}.

Canonical encodings coupled with exact supplementary evaluators also
occur in recent finite graph-reconstruction work, for example the
integer-partition model and cut-partition implementation for
diameter-six trees in \cite{Zeng2026MultiplicityArxiv}.  Here the
canonical objects are instead five-edge sets under $\Aut(T)$, and the
decisive membership test is the bundle criterion \eqref{eq:criterion}.

\section{Reproducibility and limitations}

The supplementary certificate contains the labeled edge list, the
complete $276$-line family, the eight canonical representatives and
nine-vertex witnesses, and the deterministic verifier.  A second
verifier independently reconstructs the dominating sets, all
$376{,}992$ five-edge cases, the full $48$-element automorphism group,
the orbit partition, and the witnesses.  Both computations use exact
integers; no random choice, floating-point comparison, heuristic
symmetry assumption, or untested orbit reduction enters the reported
census.

The result is deliberately finite.  It classifies minimum bondage sets
only for the explicitly displayed graph.  The literature comparison is
bounded: the closest inspected source is \cite{Srivastava2026Truncated},
whose displayed theorem and verifier provide the baseline and one
five-edge example, but not this orbit classification.  No assertion of
global bibliographic priority is needed for the theorem.

\paragraph{Computational and AI assistance.}
OpenAI Codex was used interactively in developing and checking the exact
enumeration, designing the independent verification, and editing the
manuscript.  All central enumerations were independently recomputed
using a separately written verifier.  The claims are replayable from
the explicit finite data described above, and the author takes
responsibility for the content.

\bibliographystyle{plain}
\bibliography{references}

\end{document}
